\begin{abstract}
Este trabajo analiza la dinámica estocástica de la prima de riesgo griega respecto al bono alemán, utilizando datos diarios que abarcan el periodo de la crisis de deuda soberana. Se contrastan dos enfoques metodológicos: el modelo de \textbf{Difusión Anómala} (desde la perspectiva de la Econofísica) y el modelo \textbf{GARCH(1,1) Bayesiano} (Econometría Financiera).

Los resultados empíricos revelan un comportamiento \textit{super-difusivo} con un exponente de escala $\alpha \approx 0.94$, evidenciando una fuerte persistencia y memoria a largo plazo en el mercado de renta fija. Por otro lado, la implementación Bayesiana mediante cadenas de Markov (MCMC) permite cuantificar la incertidumbre de la volatilidad condicional, capturando eficazmente los clústeres de volatilidad extremos.

El estudio concluye que, si bien el modelo de difusión ofrece un ajuste estructural superior (menor MSE) para describir la naturaleza física del fenómeno, el enfoque GARCH resulta indispensable para la gestión operativa del riesgo a corto plazo, a pesar de su elevado coste computacional.
\end{abstract}